\section{Cross-Boundary Worked Trace}
Consider three synthetic organizations observing an authentication anomaly in the same correlation bucket. Each initial signal reports a qualifying score, but the coordinator cannot determine whether rejected credentials or provider unavailability explains the failures. The trace is an illustration of the implemented catalogue and state transitions, not an additional experimental run. No raw username, address or login record is disclosed.

The first requested entry is \texttt{auth\_failure}, because it costs one unit. Suppose each source approves \texttt{high}. The source rules remain unresolved. The next cheapest entry is \texttt{failure\_kind}; its cost ties with \texttt{source\_spread}, and lexical ordering selects it first. If each source returns \texttt{unavailable}, the rules contradict the attack hypothesis. Three independent enrolled groups then produce a contradicted incident. Source spread is no longer requested. A fixed richer strategy would still attempt that remaining entry. The savings are therefore a consequence of a declared short-circuit alternative, not an inferred statistical privacy gain.

Each agent commits its approved response and debit before sending. If a response arrives but the acknowledgement is lost, the agent retransmits the stored representation under the same release identity. It does not recompute a new fine-grained answer. Different network request identifiers referring to the same snapshot, recipient and entry still resolve to that representation. If policy tightens before the retransmission, the agent can deny further delivery of the old fine value. This cannot retract the copy that the coordinator already received.

An authorized analyst snapshot may contain the approved failure kind and the attributed incident outcome. The explanation can say that the source reported service-unavailable failures and that the declared authentication rule was contradicted. It cannot say that the organization was free of every attack. The rule concerns one scenario family and observation scope, and a contradicted hypothesis does not establish the absence of unrelated threats. For an organization analyst whose scope excludes other organizations' refinements, the authorized snapshot contains a smaller evidence set than the coordinator's view.

Now suppose a source corrects its failure kind to \texttt{mixed}. The new version no longer supplies the contradiction used in the earlier three-group outcome. Correlation recomputes the incident, and its evidence-dependent explanation becomes stale. The report service checks the new digest before returning a cached report. It may render the current outcome and coverage, while declining to repeat the unsupported earlier conclusion. This illustrates how one version change can affect both a group vote and a displayed claim. The synthetic component studies do not measure the latency of this entire trace.

\begin{figure}[t]
\centering
\setlength{\unitlength}{1mm}
\begin{picture}(82,72)
\put(3,62){\framebox(72,8){\parbox{70mm}{\centering\footnotesize Versioned signal; outcome unresolved}}}
\put(3,50){\framebox(72,8){\parbox{70mm}{\centering\footnotesize Choose one cheapest untried predicate}}}
\put(3,38){\framebox(72,8){\parbox{70mm}{\centering\footnotesize Owner policy and ledger transaction}}}
\put(3,26){\framebox(72,8){\parbox{70mm}{\centering\footnotesize Approve / coarsen / deny; persist response}}}
\put(3,14){\framebox(72,8){\parbox{70mm}{\centering\footnotesize Recompute source rules and group votes}}}
\put(3,2){\framebox(72,8){\parbox{70mm}{\centering\footnotesize Stop on decision, exhaustion or expiry}}}
\put(39,62){\vector(0,-1){4}}\put(39,50){\vector(0,-1){4}}\put(39,38){\vector(0,-1){4}}\put(39,26){\vector(0,-1){4}}\put(39,14){\vector(0,-1){4}}
\put(75,18){\line(1,0){5}}\put(80,18){\line(0,1){36}}\put(80,54){\vector(-1,0){5}}
\end{picture}
\caption{Disclosure control loop. A refused request closes further requests to that source in the investigation.}
\label{fig:release-loop}
\end{figure}

\section{Accounting and Lifecycle Invariants}
For snapshot $s$, recipient $r$ and entry $e$, the organization ledger stores at most one representation, $|L_o[s,r,e]|\leq1$. Approved releases debit their catalogue cost; denied and coarsened refinements debit zero under the tested policy. Both the representation and debit belong to one local transaction. The initial signal has its own accounted release before entering the signal outbox. At commit, the per-snapshot and rolling-hour checks enforce
\begin{equation}
\sum_{x\in L_o(s)}c_x\leq B_s,\qquad
\sum_{x\in L_o[t-3600,t]}c_x\leq B_h.
\label{eq:budgets}
\end{equation}
The current deployment defaults are six units per snapshot and 2,400 units over the preceding hour. New requests, overlapping investigations and policy version changes do not clear accumulated local charges. The bound spans snapshots on one enrolled agent; it is not a cryptographic global disclosure theorem.

The coordinator and owner transactions are separate. A lost message after the owner commits can consume budget without supplying evidence in time. This leaves the incident unresolved rather than allowing an uncharged release. If the owner loses its SQLite file, the release history cannot be assumed intact; recovery needs an operator procedure before the identity resumes. Eventual retries improve delivery only while the investigation remains eligible. The paper does not prove liveness under unbounded partitions.

Evidence has a 24-hour expiry in the prototype. Retrieval masks expired evidence, while the worker persists expiry revisions. Corrections, withdrawal and permission changes alter the recipient snapshot. A report digest covers the whole authorized snapshot, so even a change to an unmentioned item can invalidate the report. This conservative approach costs regeneration; it avoids claiming a precise dependency analysis for implicit model reasoning. An audit copy is historical and must not be served as the current report without current checks.

\begin{figure}[t]
\centering
\setlength{\unitlength}{1mm}
\begin{picture}(82,60)
\put(5,50){\framebox(72,8){\parbox{70mm}{\centering\footnotesize Authorized snapshot and current digest}}}
\put(5,38){\framebox(72,8){\parbox{70mm}{\centering\footnotesize Local model proposes finite claim tuples}}}
\put(5,26){\framebox(72,8){\parbox{70mm}{\centering\footnotesize Check references, versions, kinds and values}}}
\put(5,14){\framebox(72,8){\parbox{70mm}{\centering\footnotesize Controlled report or deterministic fallback}}}
\put(5,2){\framebox(72,8){\parbox{70mm}{\centering\footnotesize Recheck scope and digest before serving}}}
\put(41,50){\vector(0,-1){4}}\put(41,38){\vector(0,-1){4}}\put(41,26){\vector(0,-1){4}}\put(41,14){\vector(0,-1){4}}
\put(5,6){\line(-1,0){4}}\put(1,6){\line(0,1){48}}\put(1,54){\vector(1,0){4}}
\end{picture}
\caption{Interpretation lifecycle. A single repair uses the same evidence; it cannot expand the recipient's authority.}
\label{fig:claim-loop}
\end{figure}

The admission invariant concerns displayed tuples: a factual claim must match the permitted set computed from the current authorized snapshot. It does not establish source truth. If the model outputs an arbitrary statement with a valid citation, membership still fails when its predicate or value is outside the finite contract. Controlled rendering preserves the admitted meaning. Allowing unrestricted paraphrasing after that check would introduce another unverified generation boundary.

\section{Closest Methods and Measurement Detail}
\begin{table*}[t]
\centering\caption{Closest mechanisms and the narrow boundary examined here.}
\label{tab:closest}\footnotesize
\begin{tabular}{p{.20\textwidth}p{.34\textwidth}p{.37\textwidth}}\hline
Method & Established mechanism & Boundary of this study\\\hline
Controlled blacklisting~\cite{freudiger} & Controlled collaborator sharing for prediction & Finite incident predicates, source-owned retry ledger and unresolved outcomes\\
SEPIA / multiparty PSI~\cite{sepia,psi} & Cryptographic aggregation or protected threshold matching & Policy-controlled categories; no equivalent cryptographic guarantee\\
LEXIS~\cite{lexis} & Security-report drafts from detector explanation evidence & Exact current recipient snapshot and finite claim admission\\
ALCE / RARR~\cite{alce,rarr} & Citation generation and evidence-based revision & A cited claim also needs exact allowed value, version and scope\\
GSAR~\cite{gsar} & Typed grounding judgement and recovery actions & Deterministic tuple membership and an explicit one-repair limit\\
Revocation study~\cite{revocation} & Empirical retrieval failures despite withdrawn status & Current snapshot invalidation before a report is reused\\\hline
\end{tabular}
\end{table*}

Table~\ref{tab:closest} compares mechanisms, not benchmark rankings. Different input assumptions and output meanings prevent importing a published accuracy figure into our finite-contract score. Cryptographic matching, detector-feature interpretation, natural-language entailment and authorization checks answer different questions. A future direct comparison needs a common task specification and must expose where a comparator cannot express the same operation.

In the eight-case pilot, withdrawal, correction and restriction each leave one unsupported claim after initial-snapshot verification. Citation prompting improves the contradicted case but worsens the corrected case relative to direct generation. Thus its lower aggregate unsupported count does not mean that every evidence condition improves. The complete method's current-state behavior is coupled to fallback throughout, so it does not isolate a successful model contribution.

Finite-claim coverage is 1.0 for template and complete, 0.979 for no-lifecycle, 0.221 for direct and 0.267 for cited output. The denominator is the programmatic allowed-claim set. Exhaustive templates are favoured by this measure; it is not analyst usefulness. A failed model call can also have zero unsupported claims because it emits none. Failure, invalid structure, raw claim count and fallback rate must accompany the support measure.

All values in the no-fault disclosure table are sums over 50 cases, while the explanation table aggregates eight constructed snapshots. These denominators must remain distinct. Neither is a population sample of independent deployments. The separate pilot cannot show that 18.2\% less disclosure preserves an analyst's useful explanation coverage, because the evaluated snapshots were not generated by those disclosure runs.

\section{Coupled Evidence States and Evaluation Gap}
The combined design has two views of incompleteness. Correlation sees whether released predicates can decide its rule, while interpretation sees what the current recipient may receive. A coordinator can have enough cross-organization evidence to record support even when an organization analyst cannot see every contributing refinement. The analyst report may state the attributed incident outcome and permitted coverage, but it must not reconstruct those hidden values. This is an authorization limit on explanation detail, not a reason to change the coordinator's outcome.

Conversely, a complete visible record does not require a model to narrate it. The fallback already constructs the permitted facts. A generator's contribution must therefore be defined in terms of a task that templates cannot satisfy adequately, such as prioritizing supported observations for a specified analyst decision. The current finite vocabulary excludes unrestricted priority judgements, and no useful-prioritization result is reported. Extending that vocabulary requires new admission rules and independent evaluation rather than permitting arbitrary commentary after verification.

A joint study must preserve the same source observations while varying disclosure policy. It should record the final rule outcome, recipient-visible snapshot, admitted claims, omitted allowed claims, fallback reason and timing for each case. Withholding a predicate may increase unresolved incidents and reduce available report detail at the same time. Those effects cannot be inferred from the two separate component totals. The denominator for claim coverage must be the snapshot actually authorized in that condition, while a separate task label measures whether the remaining information is useful for the analyst's question.

Corrections need another control. A changed snapshot can make a report stale without changing the incident outcome. An authorization change can remove visible evidence while leaving the coordinator's rule result intact. A test should exercise both and distinguish a correctly invalidated report from a missed security transition. Similarly, a model failure should leave the incident available and cause a visible template result, not enter the detector's error count. These distinctions prevent a single aggregate success measure from concealing different failure boundaries.

This is a proposed evaluation protocol. The existing eight-case pilot establishes current-state checks over constructed snapshots, and the disclosure study establishes its own policy/outcome behavior. Neither supplies the matched incident records required to estimate the joint effect. Keeping that gap explicit allows later experiments to address it without retroactively presenting component measurements as full-system validation.
